%% kapa: exam=1 round=25 year=2014 period=2 form=A subject=accounting subject_ko=회계학 seq=22 number=62
%% source: Q-Net 감정평가사 기출문제 (cst00309) artlSeq=5101207 file=extracted/1008/5101207/2014년 제25회 감정평가사 1차시험 기출문제/2교시(A형).hwp
%% conversion: kordoc 4.24.51 -> tools/gen1.py (auto); review=proofread-sonnet; flags=none
\begin{question}{62}
다음은 (주)한영의 당기 거래 내역이다. (주)한영이 무형자산으로 보고할 수 있는 상황들로만 모두 고른 것은?
\begin{framedtext}
ㄱ. 경영진이 미래효익을 기대하고 있는 고객관계 개선 관련 프로젝트에 ￦3,000 지출 \par
ㄴ. (주)부산의 장부에 자산으로 기록하지 않았던 품질향상 제조기법을 배타적 통제가능성과 함께 획득하고 ￦2,000 지급 \par
ㄷ. 기계를 ￦30,000에 구입하면서 기계제어 소프트웨어프로그램 구입을 위해 ￦3,000 추가 지급 \par
ㄹ. 신제품에 대한 광고비 ￦20,000 지급 \par
ㅁ. (주)대한의 식별가능한 순자산의 공정가치는 ￦4,000인데, (주)한영은 (주)대한의 주식 전부를 인수하기 위해 ￦7,000 지급 \par
ㅂ. (주)한영은 다른 회사로부터 실용신안권을 ￦5,000에 인수하였으며, 이 권리를 활용하여 얻은 수익 ￦10,000의 10\%인 ￦1,000을 로열티로 지급하기로 약정 \par
ㅅ. (주)세종의 장부상 금액 ￦1,000인 디자인권을 ￦5,000에 구입 \par
\end{framedtext}
\begin{choices}
  \item ㄱ, ㄴ, ㄷ
  \item ㄴ, ㄷ, ㅅ
  \item ㅁ, ㅂ, ㅅ
  \item ㄴ, ㅁ, ㅂ, ㅅ
  \item ㄱ, ㄴ, ㄷ, ㄹ, ㅁ, ㅂ, ㅅ
\end{choices}
\end{question}
\answer{4}
